% residual_loewner.tex -- round 439, Loewner/Bezout of the
% interpolant residual D0 = W_Y - S0 (reviewer DCCCIV).
% Ausgang RESIDUAL_LOEWNER_PLUS_CONTROLLED_DIAGONAL.
% Builds with pdflatex.
% Companion machine check: verify_residual_loewner.py.
% Discovery: experiments/tfpt-discovery/residual_loewner_probe.py.
% NO RH CLAIM.  Research documentation, not a theorem of RH.
\documentclass[11pt]{article}

\usepackage[a4paper,margin=2.7cm]{geometry}
\usepackage{amsmath,amssymb,amsthm}
\usepackage{microtype}
\usepackage{booktabs}
\usepackage{xcolor}
\usepackage[colorlinks=true,linkcolor=blue!50!black,urlcolor=blue!50!black]{hyperref}

\newtheorem{lemma}{Lemma}
\newtheorem{prop}{Proposition}
\theoremstyle{remark}
\newtheorem{remark}{Remark}

\newcommand{\Tfrak}{\mathfrak{T}}
\newcommand{\indm}{\operatorname{ind}_{-}}
\newcommand{\CInd}{\operatorname{CInd}_{\mathbb{R}}}
\newcommand{\YY}{Y}

\title{\bfseries The interpolant residual as Loewner plus diagonal\\[2mm]
\large $D_0=W_Y-K_{YY}^{-1}$ has displacement rank $2$;
$\mathrm{Bez}(P_Y,Q)$ does not recover $D_0$}
\author{Stefan Hamann}
\date{August 30, 2026}

\begin{document}
\maketitle

\begin{abstract}\noindent
After R431 the one negative square of the graph resolvent
lives in the interpolant residual $D_0=W_Y-S_0$ on the
holes $Y$, not in an Euler factor of the source.
This round asks whether $D_0$ is a Loewner/Bezout object,
so that $\indm(D_0)$ is zero-structure of a rational pair.
Displacement rank $\mathrm{rank}(YD_0-D_0Y)=2$ holds
exactly over $\mathbb{Q}$ and with $\sigma_3/\sigma_2\sim 10^{-14}$
on $w_9$ and two further windows.
$S_0$ is exactly the inverse of the RKHS kernel $K_{YY}$ of
$P_{<d_0}$ on the $X$-measure.
At kernel dimension $0$, Cauchy--$\pi$ additionally gives
the dressed identity
$\Delta D_0\Delta=L_{-\widetilde{m}}+\mathrm{diag}(w_Y(P_Y')^2)$
with $\widetilde{m}$ the generating function of the
$\pi_Y(X)^2$-dressed $X$-measure.
$D_0$ itself is \emph{not} a classical ones-Loewner matrix
once $|Y|\ge 3$ (the old trap: $n_Y=2$ is vacuous).
No source polynomial $Q$ with $P_Y$-denominator makes
$D_0=S^T\mathrm{Bez}(P_Y,Q)S$.
The associated rational is $m$-like with poles at $X$,
not at $Y$.
Ausgang \texttt{RESIDUAL\_LOEWNER\_PLUS\_CONTROLLED\_DIAGONAL}.
The new P1 problem is $\indm(W_Y-K_{YY}^{-1})\le 1$.
No RH claim.
\end{abstract}

\medskip\noindent
\textbf{Status.}\enspace
Lemmas~\ref{lem:disp}--\ref{lem:kinv} are proved (finite
identities).
Proposition~\ref{prop:bezout} is a named miss.
Propositions~\ref{prop:windows}--\ref{prop:census} are
finite machine checks.
Ausgang \texttt{RESIDUAL\_LOEWNER\_PLUS\_CONTROLLED\_DIAGONAL}.
No $L^*$ claim.  No $R^\dagger$ claim.  No RH claim.

\section{The residual}

Let $X$ be the occupied dual nodes of a window, $Y$ the
holes, $u^\vee$ the r356 dual weights, and $d_0=N_w-3$ the
interpolating depth of $\Tfrak_0$ (r409).
Write $W_Y=\mathrm{diag}(u^\vee|_Y)$ and $S_0$ for the
Gram of the min-norm interpolant in $P_{<d_0}$ (the
Schur short-circuit of r409).
The residual is
\begin{equation}\label{eq:D0}
D_0=W_Y-S_0,
\end{equation}
so that $I-\Tfrak_0^*\Tfrak_0$ is the $W_Y$-congruence of
$D_0$.
At kernel dimension $k_{\dim}=d_0-|Y|=0$, R431 gives the
Cauchy--$\pi$ formula
$S_0=\Delta^{-1}C\Pi W_X\Pi C^T\Delta^{-1}$ with
$\Delta=\mathrm{diag}(P_Y'(y_i))$ and $C_{ia}=1/(x_a-y_i)$.
The one negative square of the bulk sits in $D_0$.

\section{Displacement rank two}

\begin{lemma}[Displacement]\label{lem:disp}
Let $Y=\mathrm{diag}(y_i)$.
On the five-atom $k_{\dim}=0$ toy and on the six-node
$\mathbb{Q}$-window of the R431 audit,
$\mathrm{rank}(YD_0-D_0Y)=2$ exactly over $\mathbb{Q}$.
\end{lemma}

A diagonal commutes with $Y$, so
$YD_0-D_0Y=-(YS_0-S_0Y)$.
At $k_{\dim}=0$, $S_0$ is a $\Delta$-dressed Cauchy--Pick
kernel, hence Cauchy-like of displacement rank $2$.
The identity $S_0=K_{YY}^{-1}$ of Lemma~\ref{lem:kinv}
extends the same conclusion to every $k_{\dim}$: the
inverse of a Cauchy-like matrix is Cauchy-like
(Gohberg--Semencul).

\begin{prop}[Windows]\label{prop:windows}
On $w_9$ ($\lvert Y\rvert=104$, $d_0=181$, $k_{\dim}=77$),
$k_z=15$ and $k_z=18$, the singular values of $YD_0-D_0Y$
are a repeated pair $(\sigma,\sigma,0,0,\ldots)$.
The ratio $\sigma_3/\sigma_2$ is $9.42\cdot 10^{-15}$,
$1.08\cdot 10^{-14}$, $4.41\cdot 10^{-15}$ (numpy);
mpmath dps~$40$ on $w_9$ gives $9.42\cdot 10^{-15}$
with $\sigma_3=1.30\cdot 10^{-15}$.
$\Lambda$-permute keeps $\sigma_3/\sigma_2=5.71\cdot 10^{-14}$.
\end{prop}

The stop \texttt{DISPLACEMENT\_RANK\_GROWS} does not fire.

\section{$S_0$ is the inverse CD kernel}

\begin{lemma}[$S_0=K_{YY}^{-1}$]\label{lem:kinv}
Let $K_{YY}$ be the reproducing kernel of $P_{<d_0}$
for the $X$-measure, evaluated at $Y$
($K=\Phi_Y G_X^{-1}\Phi_Y^T$, $G_X$ the moment Gram).
Then $S_0=K_{YY}^{-1}$ exactly over $\mathbb{Q}$ on both
toys, and to relative residual $1.05\cdot 10^{-12}$ on
$w_9$.
Consequently
\begin{equation}\label{eq:D0K}
D_0=W_Y-K_{YY}^{-1}.
\end{equation}
\end{lemma}

This is the min-norm interpolation identity: the Gram of
coordinate interpolants in value space is the inverse of
the restricted kernel.
$K_{YY}$ is the Christoffel--Darboux kernel of the
$X$-measure, a Bezout kernel of the orthogonal pair
$(p_{d_0},p_{d_0-1})$, \emph{not} of $(P_Y,Q)$.

\section{The old trap, and the dressed Loewner}

A classical Loewner matrix of one function $f$ satisfies
$(y_i-y_j)(L_f)_{ij}=f(y_i)-f(y_j)$ (ones-normalisation).
On the five-atom toy $\lvert Y\rvert=2$ this is vacuous
(one pair): the off-diagonals determine a nodal $f$, the
polynomial extension has $L_f$ with $\indm=0$, and
$\mathrm{diag}(\Delta)=D_0-\mathrm{off}(L_f)$ carries
$\indm=1$ --- the wall-side trap of Rounds~184--191,
now on $D_0$.
On the six-node window $\lvert Y\rvert=3$ the ones-residual
is $40073273/275164890\neq 0$: $D_0$ is \emph{not}
$L_f+\mathrm{diag}(\Delta)$ in the node basis.

\begin{lemma}[Dressed Loewner at $k_{\dim}=0$]\label{lem:dress}
Let $\widetilde{m}(z)=\sum_a u^\vee(x_a)\,P_Y(x_a)^2/(x_a-z)$
and $L_{\widetilde{m}}$ its Loewner matrix at $Y$
(diagonal $\widetilde{m}'$).
At $k_{\dim}=0$,
\begin{equation}\label{eq:dress}
\Delta D_0\Delta
=L_{-\widetilde{m}}+\mathrm{diag}\bigl(w_Y(P_Y')^2\bigr)
\end{equation}
exactly over $\mathbb{Q}$.
$L_{-\widetilde{m}}$ is negative definite ($\indm=\lvert Y\rvert$);
the diagonal is positive definite.
\end{lemma}

The live pair is therefore $(L_{-\widetilde{m}},\mathrm{diag}(w_Y(P_Y')^2))$,
not a single-function Loewner of $D_0$.
The diagonal is source-pure (dual weights times $P_Y'^2$).
It does not by itself carry $\indm(D_0)$ (it is PD);
$L_{-\widetilde{m}}$ does not by itself carry it either
(it is ND).
The mixture is the P1 object.
\texttt{DIAGONAL\_SHIFT\_IS\_THE\_TARGET} does not fire
as a stop on this pair.

\section{Bezout of $(P_Y,Q)$ misses $D_0$}

The contracted B1 asks for a source polynomial $Q$ with
$D_0=S^T\mathrm{Bez}(P_Y,Q)S$ and $Q$ from the
Birkhoff/interpolation formula, never from eigendata of
$D_0$.
The Cauchy--$\pi$ rational $\widetilde{m}$ has poles at
$X$, denominator $P_X$, not $P_Y$.
The Lagrange interpolant of $\widetilde{m}$ at $Y$ is a
polynomial $Q$ of degree $<\lvert Y\rvert$ (source-first).

\begin{prop}[B1 miss]\label{prop:bezout}
For $Q$ the interpolant of $\widetilde{m}$ at $Y$, none of
the source maps $S\in\{V^{-1},\Delta^{-1}V^{-1},V\}$
($V$ the Vandermonde at $Y$) satisfies
$S^T\mathrm{Bez}(P_Y,Q)S=D_0$ on either $\mathbb{Q}$-toy.
The last Sturm remainder of $(P_Y,Q)$ on the five-atom
toy is $-7466971/10681442\neq 187/450=1-q^\dagger$ (r433).
\end{prop}

The CD pair $(p_{d_0},p_{d_0-1})$ represents $K_{YY}$,
not $D_0$.
\texttt{RATIONAL\_FUNCTION\_USES\_EIGENDATA} does not fire
($Q$ is source).
B1 does not stand, so B2 (Cauchy index of $Q/P_Y$) is not
an identity for $\indm(D_0)$.
On $w_9$, the residue count $\#\{i:\widetilde{m}(y_i)/P_Y'(y_i)<0\}=52$
is not $\indm(D_0)=1$: that CInd is a $k_{\dim}=0$ count
of holes, not P1 on the live window.
\texttt{CAUCHY\_INDEX\_RESTATEMENT} is not claimed.
The R437 edge contract (augmented $(P^\dagger,Q^\dagger)$
with last Sturm pivot $\delta=1-q^\dagger$) is not opened.

\section{Worlds}

\begin{prop}[Census]\label{prop:census}
MAIN $w_9/k_z15/k_z18$: $\indm(D_0)=1$.
$\chi_3$-$9$: $\indm=0$ (PD, world-separating).
$\chi_3$-$15$: $\indm=1$.
$\Lambda$-permute / jitter / scramble send $\indm$ to
$20/19/21$.
Displacement rank $2$ survives the kills.
\end{prop}

In zero language the hoped census
``MAIN$/\chi$ at most one non-interlacing place, permute/scramble
many'' is a statement about $\indm(W_Y-K_{YY}^{-1})$, not
about interlacing of a $P_Y$-denominator $Q$.

\section{The P1 problem}

\begin{remark}[Problem]
Prove, source-purely, $\indm(W_Y-K_{YY}^{-1})\le 1$ on the
selected sequence (equivalently at $k_{\dim}=0$: the
positive diagonal $w_Y(P_Y')^2$ lifts all but at most one
negative direction of $L_{-\widetilde{m}}$).
Sturm/Christoffel candidates: compare dual hole weights to
the Christoffel function $1/K(y_i,y_i)$ of $P_{<d_0}$ on
$X$, plus off-diagonal control of $K^{-1}$.
Gauss-node interlacing of $(p_{d_0},p_{d_0-1})$ makes $K$
PD and does not by itself bound $\indm(W_Y-K^{-1})$.
This is a finite matrix problem per window.
It is not a theorem of the Riemann hypothesis.
\end{remark}

Mincut unchanged (base~$4$ / refined~$5$).
No $L^*$ claim.  No $R^\dagger$ claim.  No RH claim.

\end{document}
